\section{Closed endpoint and rooted hinge-tree kinematics}
\label{sec:endpoint}

The endpoint \(\theta=0\) is a closed-contact endpoint, not a positive-clearance
endpoint.  The pieces are closed sets and are allowed to meet along the
catalogued contacts in Section~\ref{sec:model}.

\begin{lemma}[Closed-contact endpoint]
\label{lem:endpoint}
At \(\theta=0\), the four pieces lie in \(\T\), their volumes sum to
\(\Vol(\T)\), all six unordered pairs are exactly the catalogued contacts C0--C5,
and no two distinct pieces have strict interior overlap.  This endpoint has
zero clearance at the catalogued contacts.
\end{lemma}

\begin{proof}
By Lemma~\ref{lem:tiling}, the closed union is \(\T\) and relative interiors
are pairwise disjoint.  Thus every pairwise intersection is contained in a
proper face or edge of each piece and cannot contain a three-dimensional
interior point common to two pieces.  The contact table in Section~\ref{sec:model}
lists all six unordered piece pairs.  Since the list includes shared faces and
shared edges, the endpoint is contact-allowing and not positive-clearance.
\end{proof}

The moving model uses rooted hinge trees.  For an oriented axis from \(a\) to
\(b\), write \(u=(b-a)/\lVert b-a\rVert\).  The rotation about the affine line
through \(a,b\) by angle \(\phi\) is
\[
\Rot(a,b,\phi)(x)=Q_u(\phi)x+a-Q_u(\phi)a,
\]
where \(Q_u\) is Rodrigues' rotation matrix.  Points on the axis are fixed.

The audited representative trees are:
\begin{center}
\begin{tabular}{lll}
\toprule
Tree & Hinge graph & Signed hinge angles\\
\midrule
\treeA{} & \(P_0P_1, P_0P_2, P_1P_3\) &
\(d_{H0}=\theta,\ d_{H4}=\theta,\ d_{H7}=-\theta\)\\
\treeB{} & \(P_0P_1, P_1P_3, P_2P_3\) &
\(d_{H0}=\theta,\ d_{H7}=-\theta,\ d_{H9}=\theta\)\\
\bottomrule
\end{tabular}
\end{center}
The root is \(P_0\) in both cases.

\begin{lemma}[Well-defined rooted kinematics]
\label{lem:kinematics}
For \treeA{} and \treeB{}, any real assignment of signed angles to the selected
hinges determines a unique rigid transform \(F_i\) for each piece \(P_i\).  For
each selected hinge edge, the parent and child transforms agree on the hinge
axis.
\end{lemma}

\begin{proof}
Each selected graph has four vertices, three edges, and the rooted paths listed
in the source lock; hence it is a tree.  There is a unique path from root
\(P_0\) to every piece.  Starting with \(F_0=\mathrm{id}\), propagate along each
edge by composing the parent transform with the axis rotation about the parent-
placed hinge line.  A composition of rigid transforms is rigid, and uniqueness
of the rooted path gives uniqueness of \(F_i\).  Since \(\Rot(a,b,\phi)\) fixes
the axis line pointwise, parent and child transforms agree on the selected
hinge axis.
\end{proof}

At \(\theta=0\), every signed angle is zero, every axis rotation is the identity,
and all piece transforms are identity.  Thus both representative trees have the
same closed-contact endpoint from Lemma~\ref{lem:endpoint}.

\begin{certprop}[Kinematic source-lock]
The endpoint and kinematic statements are source-locked to
\path{certified/source_docs/S4_LEMMA_02_CLOSED_ENDPOINT.md} and
\path{certified/source_docs/S4_LEMMA_03_KINEMATICS_AND_SIGNS.md}.  These files
also record the exact hinge IDs \texttt{H0\_A\_M\_AB}, \texttt{H4\_C\_M\_CD},
\texttt{H7\_D\_M\_CD}, and \texttt{H9\_B\_M\_AB}.
\end{certprop}
